\chapter{Parameters and configurability}
\label{ch:param}

\section{Build parameters (synthesis-time)}
The hardware architecture is fixed at synthesis through the following Verilog
parameters. They determine the datapath contained in the bitstream and its capacity
\emph{ceiling}.

\begin{tabularx}{\textwidth}{L{3.0cm} C{1.8cm} Y}
\toprule
\rowh \thd{Parameter} & \thd{Default} & \thd{Meaning} \\
\midrule
\code{DATA\_WIDTH}   & 8   & Data width (INT8). \\
\rowa \code{ACC\_WIDTH} & 32 & Accumulator width (INT32). \\
\code{N\_INPUTS}     & 32 / 256 & Maximum number of inputs per neuron (benchmark baseline: 256). \\
\rowa \code{N\_NEURONS} & 1 / 4 & Maximum number of neurons per layer. \\
\code{PARALLEL}      & 8   & Simultaneous hardware MACs per neuron; must divide \code{N\_INPUTS} and should be a power of two. \\
\rowa \code{N\_LAYERS} & 4 & Maximum number of layers chainable by \code{layer\_sequencer}. \\
\code{ADDR\_WIDTH}   & 23  & Byte-address width (8~MB). \\
\rowa \code{MEM\_DATA\_WIDTH} & 16 & Width of the physical PSRAM data bus. \\
\code{CLK\_FREQ\_MHZ} & 80 & Frequency used in the PSRAM timing formulas (must be aligned to the real oscillator). \\
\bottomrule
\end{tabularx}

\begin{fnwarn}[\texttt{N\_INPUTS} \% \texttt{PARALLEL} constraint]
\code{PARALLEL} must divide \code{N\_INPUTS} exactly, otherwise the elaboration guard
fires (§\ref{ch:datapath}). The same constraint applies at runtime to
\code{n\_inputs\_real}.
\end{fnwarn}

\section{Runtime network width}
A single bitstream serves any topology \emph{up to} the build maximum. The actual width
of each execution is a separate value, set by the host:
\begin{itemize}
\item \code{n\_inputs\_real} --- inputs actually used in this execution (must be a
multiple of \code{PARALLEL});
\item \code{n\_neurons\_real} --- neurons actually computed in this execution.
\end{itemize}
Both default to the build maximum, so any caller that leaves them unconnected processes
the full width as before the ports were introduced.

\begin{fnnote}[Real early termination]
This is not mere address bookkeeping: the two values directly bound the hardware loops
(X/W reads of \code{neuron\_memory}, MAC group count of \code{neuron\_parallel} and the
length of the ping-pong copy for \code{RUN\_NETWORK}). A narrower layer actually
\emph{computes} and \emph{copies} faster and does not require zero-padding of the RAM
for the unused tail: data beyond \code{n\_inputs\_real}/\code{n\_neurons\_real} is never
read.
\end{fnnote}

This lets a network taper within a single chained execution, for example
$256\to64\to16\to4$, with each layer declaring its own actual width in the descriptor
table (ch.~\ref{ch:seq}).

\subsection{Measured savings}
Early termination was measured end-to-end:
\begin{tabularx}{\textwidth}{L{5.5cm} C{3.0cm} Y}
\toprule
\rowh \thd{Test} & \thd{Cycles} & \thd{Comparison} \\
\midrule
\code{neuron\_parallel\_tb.v} (T7) & 3 vs 6 & reduced vs full, with ``garbage'' data in the skipped lanes (proof that they are not read). \\
\rowa \code{neuron\_memory\_tb.v} (T5) & 209 vs 788 & 8-of-32 vs full 32, through the real PSRAM stack. \\
\bottomrule
\end{tabularx}

\section{Characterized configurations}
Some combinations validated in simulation and/or synthesis:
\begin{tabularx}{\textwidth}{C{2.0cm} C{2.0cm} C{2.0cm} Y}
\toprule
\rowh \thd{N\_INPUTS} & \thd{N\_NEURONS} & \thd{PARALLEL} & \thd{Notes} \\
\midrule
32 & 4 & 8 & First functional parametric test (Phase~1). \\
\rowa 256 & 4 & 2/4/8/16 & Datapath benchmark sweep (Phase~7). \\
32 & 1..3 & 8 & Single/multi-neuron memory integration (Phase~3). \\
\rowa 4 & 4 & 2 & End-to-end 2-layer \code{RUN\_NETWORK} test over real SPI. \\
\bottomrule
\end{tabularx}

\section{Build versus runtime summary}
\begin{center}
\begin{tikzpicture}[font=\footnotesize,node distance=6mm]
 \node[fnblockD,minimum width=54mm,minimum height=15mm](b){\textbf{BUILD (synthesis)}\\[2pt]
   {\scriptsize N\_INPUTS, N\_NEURONS, N\_LAYERS,}\\{\scriptsize PARALLEL, DATA\_WIDTH, ACC\_WIDTH}\\{\scriptsize $\Rightarrow$ machine ceiling}};
 \node[fnblockT,right=16mm of b,minimum width=54mm,minimum height=15mm](r){\textbf{RUNTIME (host, SPI)}\\[2pt]
   {\scriptsize n\_inputs\_real, n\_neurons\_real,}\\{\scriptsize activation, num\_layers, weights/bias}\\{\scriptsize $\Rightarrow$ actual network}};
 \draw[fnbus] (b) -- node[fnlbl,above]{$\le$} (r);
\end{tikzpicture}
\end{center}
